\documentclass[memoire.tex]{subfiles}
\begin{document}

\section*{Notations}
\addcontentsline{toc}{section}{Notations}

\small
\renewcommand{\arraystretch}{1.12}

\noindent\textbf{Cadre général.}\quad
Sauf mention contraire, $X,Y,Z$ désignent des espaces de Banach réels,
$\Lambda$ un espace de paramètres, $x\in X$ une variable d'état et
$\lambda\in\Lambda$ un paramètre. On considère typiquement une application
$F:X\times\Lambda\rightarrow Y$ et un zéro de référence $(x_0,\lambda_0)$.

\medskip
\noindent\textbf{Calcul différentiel et opérateurs}
\smallskip

\begin{tabular}{@{}p{0.27\textwidth}p{0.68\textwidth}@{}}
$L(X,Y)$
& Espace des applications linéaires de $X$ dans $Y$.\\
$\mathcal L(X,Y)$
& Espace des applications linéaires continues de $X$ dans $Y$.\\
$DF(x)$
& Dérivée de Fréchet de $F$ au point $x$.\\
$D_iF$, $D^2_{ij}F$
& Dérivées partielles de $F$ par rapport à la $i$-ième variable, et dérivées partielles secondes/mixes.\\
$X^*$, $D^*$
& Dual topologique de $X$ et opérateur adjoint de $D$.\\
$\ker D$, $\operatorname{im}D$
& Noyau et image de l'opérateur $D$.\\
$\operatorname{codim}V$
& Codimension d'un sous-espace vectoriel $V$.\\
$\operatorname{Ind}D$
& Indice de Fredholm :
$\dim\ker D-\operatorname{codim}(\operatorname{im}D)$.\\
$\mathcal F_n(X,Y)$, $\mathcal F(X,Y)$
& Opérateurs de Fredholm $X\to Y$ d'indice $n$, respectivement de tout indice.\\
$\mathcal{GL}(X)$
& Groupe des automorphismes linéaires continus de $X$.\\
$A^\perp$
& Annihilateur de $A$ dans le dual ; dans un cadre hilbertien, désigne aussi le complément orthogonal.\\
$\rangle v_1,\ldots,v_k\langle_R$
& Sous-espace ou sous-module engendré par $v_1,\ldots,v_k$ sur l'anneau (ou corps) $R$.\\
\end{tabular}

\medskip
\noindent\textbf{Réduction de Lyapunov--Schmidt et bifurcation}
\smallskip

\begin{tabular}{@{}p{0.27\textwidth}p{0.68\textwidth}@{}}
$D=D_1F(x_0,\lambda_0)$
& Linéarisation de $F$ par rapport à la variable d'état au point étudié.\\
$X=\ker D\oplus X_0$
& Décomposition de l'espace d'état en noyau de $D$ et complément fermé.\\
$Y=\operatorname{im}D\oplus Y_0$
& Décomposition de l'espace d'arrivée en image de $D$ et complément fermé.\\
$\pi_1$, $\pi_2$
& Projections associées sur $\ker D$ et $Y_0$.\\
$\varphi$
& Fonction implicite donnant la composante dans $X_0$ lors de la réduction.\\
$\Phi$
& Fonction réduite (ou fonction de bifurcation) à valeurs dans $Y_0$.\\
$F(\cdot,\lambda)$
& Application $x\mapsto F(x,\lambda)$ obtenue en fixant le paramètre.\\
$\alpha_F(x_0,\lambda_0)$
& $\dim\!\bigl(\ker D_2^*\cap\ker D_1^*\bigr)$, utilisée pour les zéros de type $(\alpha,n)$.\\
\end{tabular}

\newpage

\noindent\textbf{Germes et théorie des singularités}
\smallskip

\begin{tabular}{@{}p{0.27\textwidth}p{0.68\textwidth}@{}}
$\mathcal E(S,X,Y,T)$
& Espace des germes lisses, de source $S\subset X$ et à valeurs dans $Y$, envoyant $S$ dans $T$.\\
$\mathcal E_S(N)$
& Anneau des germes de fonctions réelles définies au voisinage de $S\subset N$.\\
$\mathcal E_n^m$
& Abréviation de $\mathcal E(0,\mathbb R^n,\mathbb R^m)$.\\
$\mathfrak m_S$, $\mathfrak m_n$
& Idéal des germes scalaires s'annulant sur la source ; $\mathfrak m_n$ désigne le cas $(0,\mathbb R^n)$.\\
$\mathcal D_n$
& Groupe des germes de difféomorphismes de $(\mathbb R^n,0)$ fixant l'origine.\\
$j^k f$
& $k$-jet de $f$, c'est-à-dire son développement de Taylor tronqué à l'ordre $k$.\\
$\mathcal K$
& Groupe de contact agissant sur les germes.\\
$\mathcal B$, $\mathcal{RB}$
& Groupes d'équivalence de bifurcation, respectivement non restreinte et restreinte (paramètres fixés).\\
$T_f(\mathcal G\!\cdot\! f)$, $T_f\mathcal G$
& Espace tangent à l'orbite de $f$ sous l'action du groupe $\mathcal G$.\\
$T_f^{(e)}\mathcal G$
& Espace tangent étendu associé à l'action de $\mathcal G$.\\
$\operatorname{codim}(f,\mathcal G)$
& $\dim_{\mathbb R}\bigl(\mathcal E_n^m/T_f\mathcal G\bigr)$, lorsque cette dimension est finie.\\
$\mathcal U_d(f)$
& Ensemble des déploiements à $d$ paramètres du germe $f$.\\
$H\preccurlyeq G$
& $H$ est un facteur du déploiement $G$ (pour la relation d'équivalence considérée).\\
\end{tabular}

\medskip
\noindent\textbf{Degré, orientation et dynamique}
\smallskip

\begin{tabular}{@{}p{0.27\textwidth}p{0.68\textwidth}@{}}
$\deg_B(f,U,y)$
& Degré de Brouwer de $f$ sur $U$ relativement à $y$.\\
$\mathcal C(D)$
& Ensemble des correcteurs de l'opérateur de Fredholm d'indice nul $D$.\\
$E\sim_D E'$
& Les correcteurs $E,E'$ définissent la même orientation de $D$.\\
$\nu(D)$, $\operatorname{sgn}D$
& Orientation naturelle d'un isomorphisme et signe d'un opérateur orienté.\\
$\deg_f(F,U,y)$
& Degré orienté de Furi d'une fonction de Fredholm d'indice nul.\\
$F\pitchfork Y_0$
& $F$ est transverse au sous-espace $Y_0$.\\[0.2em]
$V^{\mathbb C}$, $D^{\mathbb C}$
& Complexifiés d'un espace vectoriel réel $V$ et d'un opérateur réel $D$.\\
$\sigma(D)$
& Spectre de $D$ (entendu via son complexifié).\\
$\mu$
& Courbe de valeurs propres perturbées au voisinage d'un équilibre critique.\\
$C_{2\pi}^{k+\alpha}(\mathbb R,E)$
& Espace de Hölder des fonctions $2\pi$-périodiques à valeurs dans $E$.\\
$\kappa$
& Paramètre de fréquence ; après remise à l'échelle, la période vaut $2\pi/\kappa$.\\
$G(x,\kappa,\lambda)$
& Dans l'étude de Hopf, $G=\kappa\,\partial_t x-F(x,\lambda)$.\\
$J_0$
& Linéarisation de $G$ au point critique :
$J_0=\kappa_0\partial_t-D_1F(0,\lambda_0)$.\\
$S_\theta x$
& Translation de phase : $(S_\theta x)(t)=x(t+\theta)$.\\
\end{tabular}

\medskip
\noindent\textit{Attention aux notations dépendant du contexte :}
$\mathcal K$ désigne le groupe de contact dans la partie sur les singularités,
tandis que $\mathcal K(\bar U)$ désigne les applications compactes dans la
partie Leray--Schauder. De même, $A^\perp$ désigne un annihilateur dans le cadre
des espaces de Banach et un complément orthogonal lorsqu'un produit scalaire
est fixé.

\end{document}
